% =====================================================================
%  Выводы по лабораторной работе.
% =====================================================================

\labpart{Выводы}

В ходе выполнения лабораторной работы реализована автоматизированная
система анализа (распознавания) речи по варианту~1 (английский язык,
научные статьи по computer science) как продолжение ЛР8 и как ещё один
способ управлять существующей информационно-поисковой системой:
голосом можно продиктовать поисковый запрос, перейти в любой раздел
приложения, озвучить открытый документ и остановить озвучивание.
Распознавание выполняет локальная модель Whisper (faster-whisper) в том
же \texttt{speech-service}, что и синтез речи ЛР8, без обращения к
облачным сервисам.

Обязательные возможности реализованы следующим образом.
\emph{Задание списка операций}: таблица \texttt{speech\_commands} и
вкладка «Команды», на которой пары «фраза $\to$ действие» добавляются,
изменяются, отключаются и удаляются без переразвёртывания системы;
начальный набор -- 13 действий на английском и русском языках.
\emph{Автоматическая реакция с уведомлением}: голосовой режим
слушает непрерывно, разбивает речь на фразы по паузам (а не на равные
интервалы времени, которые режут длинные команды пополам), находит
команду и выполняет её, показывая распознанный текст, пульсацию
микрофона в такт голосу, зелёную волну и уведомление о сработавшей
команде. \emph{Выбор естественного языка}: английский или русский,
причём выбор определяет не только язык распознавания, но и то, какие
команды считаются действующими.

Измерения на синтезированной речи показали для основного языка
варианта (английского) 92{,}3\,\% верно найденных команд и WER 2{,}6\,\%
на предложениях из предметной области; шум с отношением сигнал/шум до
5~дБ на точность не повлиял. Русский язык, добавленный сверх
обязательного по варианту, распознаётся хуже (61{,}5\,\% команд, WER
27{,}1\,\%): причина -- не только качество распознавания, но и жёсткое
сопоставление по подстроке, не терпящее ни одного изменения окончания.
Модель \texttt{base} в потоке втрое быстрее \texttt{small}, но
почти вдвое чаще ошибается.

Проверка выявила и слабые места, описанные в разделе «Анализ
результатов и предложения по улучшению»: ложные срабатывания на
обычных предложениях, содержащих фразу команды («The search for\ldots»),
чувствительность русских команд к окончаниям и то, что все оценки
получены на синтезированной, а не живой речи. Для каждого из них
предложено решение; наиболее важные -- измерение на записанной речи и
нечёткое сопоставление команд.

В процессе проверки также найден и исправлен дефект, который
проявлялся как «ошибки через некоторое время» на микрофонах поиска и
проверки: отложенный таймер остановки предыдущего сеанса закрывал
только что начавшийся, если новая запись запускалась в течение 22~с
после остановки прежней. Теперь таймеры привязаны к своему сеансу и
отменяются при запуске нового, а сами микрофоны поиска и проверки
слушают непрерывно, как и голосовой режим, пока пользователь не
остановит их сам.
